\documentclass[11pt]{article}

\usepackage[a4paper,margin=1in]{geometry}
\usepackage{hyperref}
\renewcommand{\UrlFont}{\footnotesize}
\usepackage{enumitem}
\usepackage{titlesec}

\titleformat{\section}{\large\bfseries}{\thesection}{1em}{}
\setlist[itemize]{noitemsep, topsep=4pt}
\setlist[enumerate]{noitemsep, topsep=4pt}

\title{KCPC-CS111 \\ Problem Sheet \\ Graphs II: Dijkstra and 0--1 BFS}
\author{}
\date{}

\begin{document}

\maketitle

\noindent
Developed by the KCPC Internal Committee, this problem sheet accompanies the \textit{Graphs II} workshop.  
Problems are divided into three categories:
\begin{itemize}
    \item Integrated examples (discussed live)
    \item Core practice set (to be attempted during the workshop)
    \item Extended homework
\end{itemize}

\hrule
\vspace{1em}

% ----------------------------------------------------
\section*{I. Integrated Into Slides (Live Walkthrough)}

\begin{enumerate}
    \item \textbf{GeeksForGeeks - Implementing Dijkstra (Adjacency Matrix)}  
    
    \vspace{0.3em}
    {\url{https://www.geeksforgeeks.org/problems/implementing-dijkstra-set-1-adjacency-matrix/1}}

    \vspace{0.8em}

    \item \textbf{LeetCode 743 - Network Delay Time}  

    \vspace{0.3em}
    {\url{https://leetcode.com/problems/network-delay-time}}

        \vspace{0.8em}

    \item \textbf{LeetCode 1368 - Minimum Cost to Make at Least One Valid Path in a Grid}  

    \vspace{0.3em}
    {\url{https://leetcode.com/problems/minimum-cost-to-make-at-least-one-valid-path-in-a-grid}}

\end{enumerate}

\vspace{1em}
\hrule
\vspace{1em}

% ----------------------------------------------------
\section*{II. Core Workshop Practice Set}

These problems should be attempted during the workshop.  
They require careful modelling and correct algorithm selection.

For any questions that were not completed during the workshop, please attempt them independently afterwards.

\begin{enumerate}

    \item \textbf{LeetCode 1514 - Path With Maximum Probability}

    (Dijkstra variant with multiplicative weights)

    \vspace{0.3em}
    {\url{https://leetcode.com/problems/path-with-maximum-probability}}

    \vspace{0.8em}

    \item \textbf{LeetCode 1631 - Path With Minimum Effort}

    (Dijkstra + monotonicity reasoning)

    \vspace{0.3em}
    {\url{https://leetcode.com/problems/path-with-minimum-effort}}

    \vspace{0.8em}

    \item \textbf{LeetCode 2290 - Minimum Obstacle Removal to Reach Corner}

    (0--1 BFS application)

    \vspace{0.3em}
    {\url{https://leetcode.com/problems/minimum-obstacle-removal-to-reach-corner}}

    \vspace{0.8em}

    \item \textbf{LeetCode 542 - 01 Matrix}

    (Multi-source BFS; binary distance propagation)

    \vspace{0.3em}
    {\url{https://leetcode.com/problems/01-matrix}}

    \vspace{0.8em}

    \item \textbf{AtCoder ABC 176 D - Wizard in Maze}

    (Classic 0--1 BFS modelling)

    \vspace{0.3em}
    {\url{https://atcoder.jp/contests/abc176/tasks/abc176_d}}

    \vspace{0.8em}

    \item \textbf{Project Euler 81}

    (Weighted grid shortest path; DP or Dijkstra)

    \vspace{0.3em}
    {\url{https://projecteuler.net/problem=81}}

    \vspace{0.8em}

    \item \textbf{Project Euler 82}

    (Weighted grid shortest path; Dijkstra)

    \vspace{0.3em}
    {\url{https://projecteuler.net/problem=82}}

    \vspace{0.8em}

    \item \textbf{Project Euler 83}

    (Weighted grid shortest path; Dijkstra, 4-direction movement)

    \vspace{0.3em}
    {\url{https://projecteuler.net/problem=83}}

\end{enumerate}

\vspace{1em}
\hrule
\vspace{1em}

% ----------------------------------------------------
\section*{III. Extended Homework}

\textbf{Important.} The following problems are genuinely extremely difficult.  
They should only be attempted after completing all previously listed problems and assigned supplementary reading.  
They are intended as stretch challenges for those who wish to explore advanced modelling and optimisation techniques.

\medskip

\textbf{Matthew's Note.} I hope you enjoy working through these problems.  
If you have any questions, please feel free to contact me or the Director of Competitive Programmes (Marcell).

\vspace{0.8em}

\begin{enumerate}

    \item \textbf{LeetCode 1928 - Minimum Cost to Reach Destination in Time}

    (State expansion + constrained shortest path)

    \vspace{0.3em}
    {\small\url{https://leetcode.com/problems/minimum-cost-to-reach-destination-in-time}}

    \vspace{0.8em}

    \item \textbf{Codeforces 786B - Legacy}

    (Segment tree graph construction + Dijkstra)

    \vspace{0.3em}
    {\small\url{https://codeforces.com/problemset/problem/786/B}}

    \vspace{0.8em}

    \item \textbf{Codeforces 715B - Complete The Graph}

    (Constructive shortest path + repeated Dijkstra reasoning)

    \vspace{0.3em}
    {\small\url{https://codeforces.com/problemset/problem/715/B}}

    \vspace{0.8em}

    \item \textbf{AtCoder ABC 213 E - Stronger Takahashi}

    (Advanced 0--1 BFS modelling; non-trivial movement rules)

    \vspace{0.3em}
    {\small\url{https://atcoder.jp/contests/abc213/tasks/abc213_e}}

\end{enumerate}

\vspace{1em}

\hrule
\vspace{1em}

\noindent
\textbf{Guidance:}
\begin{itemize}
    \item Identify whether weights are uniform, binary, or general non-negative.
    \item Choose between BFS, 0--1 BFS, or Dijkstra accordingly.
    \item Focus on modelling transitions correctly before coding.
\end{itemize}

\end{document}
